% Generated draft: requires mathematical and source review.
\documentclass[11pt]{article}
\usepackage[margin=25mm]{geometry}
\usepackage{fontspec}
\setmainfont{Latin Modern Roman}
\usepackage{amsmath,amssymb,mathtools,unicode-math}
\setmathfont{Latin Modern Math}
\usepackage{enumitem,needspace,xurl,hyperref}
\setlength{\parindent}{0pt}
\setlength{\parskip}{5pt}
\setlength{\emergencystretch}{3em}
\newcommand{\N}{\mathbb{N}}
\newcommand{\Z}{\mathbb{Z}}
\newcommand{\Q}{\mathbb{Q}}
\newcommand{\R}{\mathbb{R}}
\newcommand{\C}{\mathbb{C}}
\newcommand{\F}{\mathbb{F}}
\newcommand{\A}{\mathbb{A}}
\newcommand{\PP}{\mathbb{P}}
\newcommand{\E}{\mathbb{E}}
\newcommand{\eps}{\varepsilon}
\newcommand{\dd}{\,\mathrm{d}}
\begin{document}
\section*{UnsolvedMath ID: 30003944}
\textbf{Source problem code:} catalogue entry
\subsection{Definitions and mathematical statement}
Let \(g\) denote the genus and \(k\) a grading degree. The supplied record names two graded Lie algebras \(\mathfrak t_g\) and \(\mathfrak m_g\), describing them respectively as the Torelli and symplectic-derivation graded Lie algebras. Write \(\mathfrak t_g(k)\) and \(\mathfrak m_g(k)\) for their homogeneous components of degree \(k\). The record does not specify the coefficient ring or field, the constructions or filtrations defining these algebras, or a comparison homomorphism. In particular, the designation of \(\mathfrak m_g\) does not determine whether the full symplectic-derivation algebra or a particular subalgebra or image is intended.

The supplied conjectural comparison is
\[
\mathfrak t_g(k)\cong\mathfrak m_g(k)
\qquad\text{in the stable genus range, for every degree }k\ne 2.
\]
For clarity, let \(I\) denote the unspecified set of grading degrees and \(G\) the unspecified set of admissible genera. If ``in the stable genus range'' means sufficiently large genus separately for each degree, the quantifier pattern would be
\[
\forall k\in I\setminus\{2\}\;\exists g_0(k)\in G\;
\forall g\in G,\quad
 g\ge g_0(k)\Longrightarrow
 \mathfrak t_g(k)\cong\mathfrak m_g(k).
\]
This is a conditional interpretation of the stable-range wording, not a supplied stability bound or a verified choice of quantifiers. The record does not distinguish this interpretation from a uniform genus threshold or an explicitly prescribed range relating \(g\) and \(k\).

The meaning of \(\cong\) also requires clarification: the statement uses an isomorphism symbol, whereas the supplied literature assessment refers to ``the map'' without defining that map. It is therefore unresolved whether the intended assertion concerns an abstract isomorphism of components or the isomorphism property of a specific comparison map, and which additional structures must be preserved. A homogeneous component of a graded Lie algebra need not itself be a Lie subalgebra. No conclusion in degree \(2\) is asserted. These missing specifications prevent a uniquely determined, fully self-contained mathematical formulation.
\subsection{Short English statement}
The supplied conjecture asks whether the degree-\(k\) components of two Lie algebras associated with Torelli theory agree in sufficiently large genus for every \(k\ne2\). The precise algebras, coefficients, comparison map, and stable range require clarification. The supplied open status and low-degree claims are dated curation, not independently verified conclusions.
\subsection{Sources}
\begin{itemize}
\item Supplied local UnsolvedMath registry record, permanent ID \(30003944\), title ``Stable-Range Comparison of Torelli Lie Algebras,'' source problem code \texttt{catalogue entry}, in the Oberwolfach Reports Open Problems collection. The record gives proposed year \(2018\) and supplies no individual proposer.
\item Supplied report citation: \emph{New Trends in Teichm\"uller Theory and Mapping Class Groups} (2019), \emph{Oberwolfach Reports}, pp.~2475--2534, DOI \texttt{10.4171/owr/2018/40}. \url{https://doi.org/10.4171/owr/2018/40}. The background describes these reports as documenting mathematical workshops and their open-problem discussions. The literature-source label for this link is ``Oberwolfach Report 40, Torelli problem.'' No exact problem page or speaker attribution is supplied.
\item Supplied literature source: ``Morita--Sakasai--Suzuki, first Galois obstruction,'' \url{https://arxiv.org/abs/2608.17673}. This reproduces the supplied attribution and link; the record gives no fuller bibliographic citation.
\item Literature assessment checked at \texttt{2026-08-21}: this date and assessment are supplied dated curation, not independently verified. The record labels the problem open, states that the comparison map is known to be an isomorphism in low degrees including \(k=3,4,5,6\), attributes the conjecture for every \(k\ne2\) to the report, and states that later Johnson-cokernel work does not prove the full conjecture. The supplied background calls this literature triage derived from OWR curation and explicitly requests independent verification of time-sensitive status and cited claims.
\end{itemize}
\end{document}
